\documentclass[12pt]{article}

\usepackage{amsmath}
\usepackage{booktabs}
\usepackage{listings}
\usepackage{xcolor}
\usepackage{enumitem}

\lstset{
  basicstyle=\ttfamily\small,
  breaklines=true,
  frame=single,
  backgroundcolor=\color{gray!10}
}

\title{RelayGuard ML System Documentation (Engineer Guide)}
\author{RelayGuard Team}
\date{}

\begin{document}

\maketitle
\tableofcontents
\newpage

%--------------------------------------------------

\section{System Overview}

\[
\text{Sensor Data} \rightarrow \text{Physics Model (HI)} \rightarrow \text{ML Model} \rightarrow \text{RUL}
\]

\textbf{Explanation:}

The system does \textbf{not} directly predict failure from raw sensor data. Instead:

\begin{enumerate}
  \item Sensor data is converted into a \textbf{Health Index (HI)} using physics models.
  \item The ML model learns patterns from HI and engineered features.
  \item The final output is the \textbf{Remaining Useful Life (RUL)}.
\end{enumerate}

%--------------------------------------------------

\section{Input Features (What We Measure)}

\begin{table}[h]
\centering
\begin{tabular}{lllll}
\toprule
Parameter & Symbol & Meaning & Used In \\
\midrule
Current & $I$ & Load through relay & Arc damage model \\
Voltage & $V$ & Switching voltage & Arc damage model \\
Temperature & $T$ & Relay heat & Thermal model \\
Cycles & $N$ & Number of operations & Mechanical wear \\
Resistance & $R$ & Contact degradation & Direct wear indicator \\
Temp Rate & $dT/dt$ & Heating speed & ML anomaly detection \\
Peak Current & $I_{\text{peak}}$ & Current spikes & ML stress feature \\
Frequency & $f$ & Switching rate & ML usage intensity \\
\bottomrule
\end{tabular}
\end{table}

%--------------------------------------------------

\section{Physics-Based Modeling (Core Logic)}

\subsection{Arc Damage Model}

\[
E = I^2 \cdot V \cdot t
\]

\textbf{Why used:}
\begin{itemize}
  \item Models \textbf{contact erosion during switching}.
  \item Damage increases rapidly with current.
\end{itemize}

\[
W_{total} = \sum (E \cdot K_T)
\]

\textbf{Meaning:}
\begin{itemize}
  \item Total accumulated damage over the relay's lifetime.
  \item Used as a long-term degradation indicator.
\end{itemize}

%--------------------------------------------------

\subsection{Thermal Model (Arrhenius)}

\[
K_T = e^{\left(\frac{E_a}{k}\left(\frac{1}{T_{ref}} - \frac{1}{T}\right)\right)}
\]

\textbf{Why used:}
\begin{itemize}
  \item Temperature accelerates aging of contact materials.
  \item High temperature = faster failure rate.
\end{itemize}

%--------------------------------------------------

\subsection{Resistance Model}

\[
R = R_0 + kN^\alpha
\]

\textbf{Why used:}
\begin{itemize}
  \item Contact resistance increases with each switching cycle.
  \item Provides a direct, measurable degradation signal.
\end{itemize}

%--------------------------------------------------

\subsection{Composite Health Index (HI)}

\[
HI =
0.40 \cdot HI_{damage}
+ 0.25 \cdot HI_{thermal}
+ 0.25 \cdot HI_{resistance}
+ 0.10 \cdot HI_{cycles}
\]

\textbf{Why used:}
\begin{itemize}
  \item Combines all degradation mechanisms into a single value.
  \item Gives a \textbf{single health score (0--100\%)} for easy monitoring.
\end{itemize}

%--------------------------------------------------

\section{Where Each Parameter is Used}

\begin{table}[h]
\centering
\begin{tabular}{lll}
\toprule
Parameter & Used In & Purpose \\
\midrule
$I$ & Damage model & Arc energy \\
$V$ & Damage model & Arc intensity \\
$T$ & Thermal model & Aging acceleration \\
$N$ & Cycle model & Mechanical wear \\
$R$ & Resistance model & Direct degradation \\
$dT/dt$ & ML only & Detect abnormal heating \\
$I_{\text{peak}}$ & ML only & Capture spikes \\
$f$ & ML only & Usage rate \\
\bottomrule
\end{tabular}
\end{table}

%--------------------------------------------------

\section{ML Training Strategy}

\subsection{Key Concept}

We do \textbf{NOT} train on raw sensor data directly. Instead:

\[
\textbf{Processed Features} \rightarrow \textbf{Future HI or RUL}
\]

\textbf{Why:}
\begin{itemize}
  \item The physics model provides \textbf{structured, meaningful data}.
  \item ML focuses on learning \textbf{patterns over time}.
  \item This approach improves accuracy and generalization.
\end{itemize}

%--------------------------------------------------

\section{Feature Engineering}

Final features used for ML training:

\begin{itemize}
  \item Current ($I$)
  \item Voltage ($V$)
  \item Temperature ($T$)
  \item Cycles ($N$)
  \item Resistance ($R$)
  \item Temperature rate ($dT/dt$)
  \item Peak current ($I_{\text{peak}}$)
  \item Frequency ($f$)
  \item Cumulative damage ($W_{total}$)
\end{itemize}

\textbf{Purpose:} Capture both \textbf{instant stress} (peaks, rates) and \textbf{long-term degradation} (cumulative damage, resistance).

%--------------------------------------------------

\section{Dataset Structure}

\begin{lstlisting}[language=Python]
features = [
    "current", "voltage", "temp",
    "cycles", "resistance",
    "dT_dt", "peak_current", "frequency",
    "damage"
]

target = "HI_future"
\end{lstlisting}

\textbf{Explanation:}
\begin{itemize}
  \item \textbf{Input} = current system state (sensor readings + derived features).
  \item \textbf{Output} = future degradation (HI at next time step).
  \item Enables prediction \textbf{before} failure actually occurs.
\end{itemize}

%--------------------------------------------------

\section{Model Training}

\subsection{Option 1: XGBoost (Recommended for Round 1)}

\begin{lstlisting}[language=Python]
from xgboost import XGBRegressor

model = XGBRegressor(
    n_estimators=100,
    max_depth=5,
    learning_rate=0.1
)

model.fit(X_train, y_train)
\end{lstlisting}

\textbf{Why XGBoost:}
\begin{itemize}
  \item Works well with tabular/structured data.
  \item Fast to train and highly accurate.
  \item Does not require a huge dataset to perform well.
\end{itemize}

%--------------------------------------------------

\subsection{Option 2: LSTM (Advanced)}

\begin{itemize}
  \item \textbf{Input:} time-series window (e.g., last 50 readings)
  \item \textbf{Output:} future HI or RUL
\end{itemize}

\textbf{Why LSTM:}
\begin{itemize}
  \item Captures temporal patterns and trends in sensor data.
  \item Better suited for long-term RUL prediction.
\end{itemize}

%--------------------------------------------------

\section{Prediction Pipeline}

\[
\text{Sensor Data} \rightarrow HI \rightarrow \text{ML Model} \rightarrow \text{RUL}
\]

\textbf{Steps:}
\begin{enumerate}
  \item Read live sensor data.
  \item Compute Health Index (HI) using physics models.
  \item Generate engineered features from HI and sensors.
  \item Feed features into the ML model to predict future HI/RUL.
\end{enumerate}

%--------------------------------------------------

\section{Alert System}

\begin{table}[h]
\centering
\begin{tabular}{ll}
\toprule
HI (\%) & Status \\
\midrule
$> 70$ & Healthy \\
$40$--$70$ & Warning \\
$20$--$40$ & Critical \\
$< 20$ & Failure Imminent \\
\bottomrule
\end{tabular}
\end{table}

\textbf{Usage:}
\begin{itemize}
  \item Drives real-time dashboard alerts for operators.
  \item Enables proactive maintenance scheduling.
  \item Can trigger automatic relay switching logic.
\end{itemize}

%--------------------------------------------------

\section{Full System Pipeline}

\begin{enumerate}
  \item \textbf{Sensors} --- collect current, voltage, temperature, resistance, cycles
  \item \textbf{Physics Model (HI)} --- compute arc damage, thermal aging, resistance wear
  \item \textbf{Feature Engineering} --- combine raw + derived features
  \item \textbf{ML Model} --- learn degradation patterns
  \item \textbf{Prediction (RUL)} --- estimate remaining useful life
  \item \textbf{Alerts} --- trigger warnings, maintenance actions
\end{enumerate}

%--------------------------------------------------

\section{Key Insight}

\[
\text{Physics Model} = \text{Ground Truth}
\]

\[
\text{ML Model} = \text{Pattern Learning}
\]

\textbf{Interpretation:}
\begin{itemize}
  \item Physics ensures \textbf{correctness} and domain knowledge.
  \item ML improves \textbf{prediction accuracy} from historical patterns.
  \item Together $\rightarrow$ reliable, explainable predictive maintenance.
\end{itemize}

%--------------------------------------------------

\section{Final Understanding}

\begin{itemize}
  \item Physics converts raw signals into \textbf{meaningful degradation indicators}.
  \item ML learns how degradation \textbf{evolves over time}.
  \item The combined system predicts relay failure \textbf{before it happens}, enabling proactive maintenance and preventing unplanned downtime.
\end{itemize}

\end{document}
